\chapter{留白、行距与导言区}

我刚开始写数学笔记时，总觉得公式能编译出来就算完成。等到复习，才发现有些页根本不想看：段落挤，分式挤，连公式编号也挤。我开始认真调整页边距和行距，后来把这些设置固定下来，写新章节时便不用从头折腾。

这份模板沿用高数、线代、概率论三套笔记的深蓝与青绿配色。定理、方法和计算提醒各有固定样式；公式与图表自动编号，PDF 中的引用可以点击。封面的图片和文字分开处理，改书名时不用重做整张图。

导言区已经加载第二至四章示例所需的数学、图表宏包，也定义了概念框和向量命令。迁移这些示例时，取文档起止命令之间的内容，按书中的章节层级调整标题。算法、子图和文献示例用到的额外宏包，则按各章的说明加入导言区。

这是学习用书和复习笔记的模板。参赛稿有自己的格式要求，第七章会单独处理摘要页。

\section{先编译一份完整文档}

把代码存为 main.tex，编码设为 UTF-8，选择 XeLaTeX，连续编译两遍。默认使用已有的中文字体；缺少 Windows 字体时会回退到 Fandol。需要的宏包应已安装，不用开启 shell escape。

封面图不是必需文件。有图就把它命名为 cover\_banner.png，放在 main.tex 旁边；没有图也能编译。我建议你先确认正文、目录和引用都正常，再换自己的封面。

本章末尾给出了\hyperref[sec:golden-template-source]{完整模板源码}。本书封面使用同一套设计，下图展示示例的正文页。
\examplepreview{previews/ch01_01_example-p3.png}

\section{geometry：把页边距算清楚}

这套参数采用 A4 纸，内侧 26 mm，外侧 23 mm。我把靠书脊的一侧多留了 3 mm，方便翻页和薄册装订。双面模式下，奇数页与偶数页会自动交换内外侧，页眉也随之镜像。

{\centering\begin{tabular}{@{}>{\raggedright\arraybackslash}p{.21\linewidth}>{\raggedright\arraybackslash}p{.29\linewidth}>{\raggedright\arraybackslash}p{.43\linewidth}@{}}
\toprule
参数 & 取值 & 含义 \\
\midrule
inner / outer & 26 / 23 mm & 靠书脊与纸张外侧的边距 \\
top / bottom & 各 25 mm & 含页眉页脚区域到纸张上下边缘的距离 \\
includeheadfoot & 开启 & 计算总版心时计入页眉、页眉间隔与页脚距离 \\
headheight & 16 pt & 页眉预留高度 \\
headsep & 6 mm & 页眉与正文之间的距离 \\
footskip & 10 mm & 正文区域底边到页脚基线的距离 \\
\bottomrule
\end{tabular}\par}

includeheadfoot 最容易被忽略。开启后，上下 25 mm 并不等于正文离纸边各 25 mm。A4 的高度是 297 mm，扣掉上下边距得到 247 mm，这部分还要容纳页眉与页脚所需的空间，所以真正的正文高度会更小。

正文宽度是 210 mm 减去 26 mm 和 23 mm，等于 161 mm。碰到特别长的公式，我会先尝试按推导步骤拆行。

做单面活页笔记时，可以把 twoside 改成 oneside，并同步调整页眉位置。厚册胶装则应根据纸张、页数和装订方式重新评估内边距。这组数值是我的起点，不是通用的“黄金比例”；模板也没有另外叠加 bindingoffset，免得装订空间加了两次。

\section{setspace：行距、段距分开调}

中文笔画密，公式又有分式和上下标，行距太小会明显影响阅读。我把 ctex 的初始倍率设为 1，再用 setspace 统一设成 1.35，免得几处行距命令相互覆盖。

\textbf{1.35 乘的是当前字号的标准基线距离。} 这套 11 pt 文档的标准基线为 13.6 pt，调整后约为 18.36 pt。遇到较高的行内公式，实际基线距离还可能增加。

我建议你拿两页真实内容试排，别只看一行测试文字。文字多时可以稍收紧，分式和推导多时再适当放宽。

段距单独设置为约 3 pt，并保留两字首行缩进；表格行高交给 arraystretch，多行公式的额外间隔交给 jot。不要指望一个行距倍率把所有间隔都调好。

我还开启了 raggedbottom。内容不足一页时，允许页底留白，尽量不为了齐底而把段落间距拉大。这里的留白由版式处理，不靠反复敲空行来凑。

\section{hyperref：引用、书签与文件信息}

当笔记请读者回看式~(7.23) 时，超链接可以省去来回翻页。hyperref 能让读者点击目录或公式编号直接跳转，也能打开正文里的网址。

colorlinks 用文字颜色标识链接；linkcolor 管内部引用，urlcolor 管网址，citecolor 管文献引用。我把三者统一为深青绿。本例尚未接入文献数据库，citecolor 先预留。

bookmarksnumbered 为侧边栏书签显示编号，bookmarksopen 与 bookmarksopenlevel 控制默认展开方式。linktoc 设为 all，目录标题和页码就都能点击。侧边栏书签与正文目录各自独立，阅读器还可能记住自己的展开偏好。

pdftitle、pdfauthor、pdfsubject 写入 PDF 元数据，封面文字仍需单独排版。做匿名稿时，署名、页眉和元数据都要检查，只删封面上的姓名还不够。

屏幕版保留青色链接。打印时把开关改为 true，hidelinks 会隐藏颜色和边框，点击功能还在。这个开关只改链接外观，不会自动把学习模板变成竞赛稿。

我把 hyperref 放在常规数学、图表和版式宏包后面。它不必永远是最后一个宏包；例如以后加入 cleveref，就应按相应说明处理加载顺序。

\section{这套样式怎样使用}

定理框放条件与结论，方法卡放可执行步骤，踩坑框只提醒确实容易出错的地方。它们都允许跨页，不过长篇推导仍适合用普通正文。每段都装进彩框，层次反而不清楚。

我保留了三部曲中的配色和卡片结构，同时加入公式引用、图表编号和三线表。示例先交代适用条件，再给公式和解释。表中的误差数值只用于演示版式，不是实验结果。

封面命令在导言区定义，在正文开始后调用。第一次改模板，先替换书名和内容；熟悉之后再调 TikZ 坐标。我不建议你同时修改十几处设置，编译出了问题很难定位。

\section{我会先检查这些地方}

编译失败，先看编译器与编码。本书统一使用 XeLaTeX 和 UTF-8。首行的编辑器提示不一定会替你切换实际编译器，仍要检查下拉菜单。中文并非只能用 XeLaTeX 排，但整本书沿用一条编译路线，排错会省事得多。

目录或引用出现问号，先再编译一次；还没解决，再查标签拼写与重复定义。导言区也别整段叠加：原模板已加载的宏包，不用为了“更保险”再加载一遍。旧资料中的 pdftex 驱动选项尤其不能机械搬到 XeLaTeX 配置里。

找不到图片，就核对文件名、扩展名和路径。把 JPG 的后缀改成 PNG，不会改变它的文件格式。

排完一定看 PDF。长公式越界就拆行，别把字号缩成蚂蚁开会。正式交稿前，再按当届要求检查匿名信息、边距、摘要与附录。

\section*{参考资料与配置核对}

版式参考数学三部曲笔记，参数含义结合中文入门资料与宏包说明核对：

\begin{itemize}
\item \href{https://ctan.org/pkg/geometry}{geometry}：页面尺寸与边距。
\item \href{https://ctan.org/pkg/setspace}{setspace}：行距控制。
\item \href{https://ctan.org/pkg/hyperref}{hyperref}：链接、书签与元数据。
\item \href{https://ctan.org/pkg/ctex}{CTeX}：中文文档类、字体与初始行距。
\end{itemize}

\section*{完整模板源码}\label{sec:golden-template-source}
\sourceattachment{examples/ch01/01_example.tex}
\VerbatimInput[baselinestretch=1.04]{examples/ch01/01_example.tex}
